\section{Supplementary coefficient and tail estimates}
\label{sec:scientific-support}

The proof uses the coefficient enclosures and omitted-mode estimates
in Sections~\ref{tail:section} and \ref{sec:computation}.  This appendix
records an independent integer recurrence for the same finite coefficient
sums and alternative exact bounds for the same omitted-mode series.  These
arguments are supplementary: they do not add a hypothesis to the scalar
theorem or replace the analytic derivative bounds and finite retained-mode
premises used there.

\subsection{An independent recurrence for the lattice coefficients}
\label{supp:coefficient-recurrence}

Fix integers $(n,q)\ne(0,0)$ and write
\[
 A=2n+q,\qquad N=n^2+nq+q^2,\qquad
 \zeta=A+i\sqrt3q,\qquad w=\frac{\zeta}{2}.
\]
Thus $w=n+q/2+i\sqrt3q/2$, $N=|w|^2>0$, and
$|\zeta|^2=4N$.  A point of the covolume-one triangular lattice is
$aw$, where $a=\sqrt{2/\sqrt3}$.  For $m\geq0$, put
\[
 r_m(n,q)=\operatorname{Re}(\zeta^{6m}).
\]
The two numbers $u=\zeta^6$ and $v=\overline{\zeta}^{6}$ have
\[
 u+v=2r_1,\qquad uv=(4N)^6,
\]
and hence satisfy $X^2-2r_1X+(4N)^6=0$.  Multiplying this equation
by $u^{m-1}$ and $v^{m-1}$, respectively, and adding gives the integer
recurrence
\begin{equation}
 \begin{split}
 r_0&=1,\\
 r_1&=A^6-45A^4q^2+135A^2q^4-27q^6,\\
 r_{m+1}&=2r_1r_m-(4N)^6r_{m-1}\qquad(m\geq1).
 \end{split}
 \label{supp:real-recurrence}
\end{equation}
Here all the quantities on the right are integers.  This proves the
recurrence for every individual lattice point, without a numerical
approximation to a complex power.

The binomial theorem identifies its numerator exactly.  Since only even
powers of $i\sqrt3q$ have nonzero real part,
\begin{equation}
 r_m(n,q)=\sum_{j=0}^{3m}(-1)^j\binom{6m}{2j}
                 A^{6m-2j}(3q^2)^j.
 \label{supp:binomial-numerator}
\end{equation}
This is the numerator in \eqref{comp:coefficient-summand}, after writing
$q$ for its second integer coordinate.  Indeed,
$w^{-6m}=\overline{w}^{6m}/N^{6m}$, and therefore
\begin{equation}
 \operatorname{Re}(w^{-6m})
 =\frac{r_m(n,q)}{2^{6m}N^{6m}}
 =\frac{r_m(n,q)}{(64N^6)^m}.
 \label{supp:reciprocal-power}
\end{equation}
Thus the recurrence and the binomial expression agree for every point and
every mode, not only for a selected finite list of comparisons.

\paragraph{Grouping by the norm.}
For a positive integer $D$, let
\[
 \mathcal Q_D=\{(n,q)\in\Z^2:\ |n|,|q|\leq D,\ (n,q)\ne(0,0)\},
 \qquad
 R_{m,N}(D)=\sum_{\substack{(n,q)\in\mathcal Q_D\\n^2+nq+q^2=N}}
                         r_m(n,q).
\]
After forming the numerator for each point, one may combine the equal
denominators in \eqref{supp:reciprocal-power}:
\begin{equation}
 \sum_{(n,q)\in\mathcal Q_D}\operatorname{Re}(w^{-6m})
 =\sum_{N\geq1}\frac{R_{m,N}(D)}{(64N^6)^m}.
 \label{supp:norm-grouping}
\end{equation}
Only finitely many positive norms occur.  It is important that the
recurrence precede this grouping.  For fixed $D$, summing
\eqref{supp:real-recurrence} over a norm shell gives
\[
 R_{m+1,N}(D)
 =2\sum_{\substack{(n,q)\in\mathcal Q_D\\n^2+nq+q^2=N}}
       r_1(n,q)r_m(n,q)-(4N)^6R_{m-1,N}(D).
\]
In general the first sum cannot be replaced by one value of $r_1$ times
$R_{m,N}(D)$.  For example, $(7,0)$ and $(5,3)$ both have $N=49$, but
their respective values of $r_1$ are $7\,529\,536$ and $-4\,912\,064$.
One may apply the recurrence either pointwise or within groups having the
same pair $(N,r_1)$, and only then combine all terms of a common norm.

For the finite coefficient enclosures, take $D=256$ when $m=1$ and $D=48$
when $2\leq m\leq38$.  The omitted dimensionless sum has absolute value
at most
\[
 E_{m,D}=8(4/3)^{3m}\frac{D^{2-6m}}{6m-2}
\]
by \eqref{tail:coefficient-square}.  The coefficient $c_m$ is consequently
enclosed by adding $[-E_{m,D},E_{m,D}]$ to the finite sum in
\eqref{supp:norm-grouping} and then multiplying by the positive factor
$a^{-6m}/(6m)$.  The square and its omitted tail are unchanged by grouping.

\paragraph{Complex pairs and interval containment.}
The finite calculation may instead keep both real components.  Write
\[
 \zeta^6=r_1+i\sqrt3\,s_1,\qquad
 s_1=6A^5q-60A^3q^3+54Aq^5.
\]
The multiplication rule
\[
 (x+i\sqrt3y)(x'+i\sqrt3y')
 =(xx'-3yy')+i\sqrt3(xy'+yx')
\]
shows that, starting from $(x_0,y_0)=(1,0)$,
\[
 (x_{m+1},y_{m+1})
 =(x_mr_1-3y_ms_1,\ x_ms_1+y_mr_1)
\]
represents $\zeta^{6(m+1)}$ exactly.  In particular $x_m=r_m$ and its real
component obeys \eqref{supp:real-recurrence}.  This is the complex-pair
recursion used in the finite calculation, again applied per point before
norm grouping.  The half square $n>0$, or $n=0$ and $q>0$, contains one
point of each pair $(n,q),(-n,-q)$.  Because $6m$ is even, the two real
numerators agree, which justifies doubling each half-square contribution.

For completeness, grouping also explains an interval containment relation.
At fixed denominator $Q=2^{144}$, let
\[
 \mathcal I_Q(x)=
 \left[\frac{\lfloor Qx\rfloor}{Q},
       \frac{\lceil Qx\rceil}{Q}\right]
 \qquad(x\in\mathbb Q).
\]
For any finite rational list $x_1,\ldots,x_h$, the elementary inequalities
\[
 \left\lfloor\sum_jQx_j\right\rfloor\geq\sum_j\lfloor Qx_j\rfloor,
 \qquad
 \left\lceil\sum_jQx_j\right\rceil\leq\sum_j\lceil Qx_j\rceil
\]
give
\begin{equation}
 \mathcal I_Q\left(\sum_jx_j\right)
       \subseteq\sum_j\mathcal I_Q(x_j).
 \label{supp:interval-grouping}
\end{equation}
Apply this to the exact rationals in each norm shell.  Converting the shell
sum once and then summing shells is contained in the enclosure obtained by
converting every point separately.  Adding the same outward tail interval
and multiplying by the same positive scaling enclosure preserve containment,
because interval addition, multiplication, and outward rounding are monotone
under inclusion.  This statement assumes the same primitive rounding rule,
tail enclosure, and scaling enclosure on both sides; exact algebra alone
does not assert endpoint identity or containment between different choices
of those enclosures.

\subsection{Alternative exact bounds for the omitted modes}
\label{supp:tail-bounds}

For the allowed indices $k=234,240,\ldots$, split the weighted summand
$(k+1)^4w_k$ of \eqref{tail:w-definition} into four positive rationals:
\begin{align*}
 t_1(k)&=10^6\,30(k+1)^4k^2(13/20)^k,\\
 t_2(k)&=10^6(3\cdot10^6)(k+1)^4k^3(1/2)^k,\\
 t_3(k)&=10^6\,5600(k+1)^4k^5(3/5)^k,\\
 t_4(k)&=10^6(7\cdot10^{11})(k+1)^4k^2(289/500)^k.
\end{align*}
Then
\[
 \mathfrak a=\sum_{j=0}^{\infty}\sum_{i=1}^4t_i(234+6j).
\]
This is the same $\mathfrak a$ as in \eqref{tail:weighted-sum}, distinct
from the lattice scale $a$.  If
\[
 (p_i,\beta_i)=(2,13/20),(3,1/2),(5,3/5),(2,289/500),
\]
in the order $i=1,2,3,4$, the exact consecutive ratios are
\begin{equation}
 \rho_i(k)=\frac{t_i(k+6)}{t_i(k)}
 =\left(\frac{k+7}{k+1}\right)^4
  \left(\frac{k+6}{k}\right)^{p_i}\beta_i^6.
 \label{supp:component-ratios}
\end{equation}
Both rational factors decrease for $k>0$, so each $\rho_i(k)$ decreases.
The four initial ratios, with no rounding, are
\[
 \begin{array}{c|c}
 i&\rho_i(234)\\ \hline
 1&(241/235)^4(40/39)^2(13/20)^6\\
 2&(241/235)^4(40/39)^3(1/2)^6\\
 3&(241/235)^4(40/39)^5(3/5)^6\\
 4&(241/235)^4(40/39)^2(289/500)^6.
 \end{array}
\]
Here is an exact simultaneous check that all four are less than $1/10$.
The integer comparisons
\begin{align*}
 9\cdot241^4=30\,360\,623\,049
   &<30\,498\,006\,250=10\cdot235^4,\\
 7\cdot40^5=716\,800\,000
   &<721\,793\,592=8\cdot39^5,\\
 800\cdot13^6=3\,861\,447\,200
   &<4\,032\,000\,000=63\cdot20^6
\end{align*}
give, since $p_i\leq5$ and $0<\beta_i\leq13/20$,
\begin{equation}
 \rho_i(234)
 \leq(241/235)^4(40/39)^5(13/20)^6
 <\frac{10}{9}\frac87(13/20)^6<\frac1{10}
 \quad(1\leq i\leq4).
 \label{supp:initial-ratio-bound}
\end{equation}

\paragraph{Three geometric bounds.}
Set $T_0=\sum_{i=1}^4t_i(234)$.  One convenient entirely integer form of
its exact comparison is obtained by putting
\begin{align*}
 C_{234}={}&30\cdot234^2\cdot325^{234}
 +3\cdot10^6\cdot234^3\cdot250^{234}\\
 &+5600\cdot234^5\cdot300^{234}
 +7\cdot10^{11}\cdot234^2\cdot289^{234}.
\end{align*}
The common denominator $500^{234}$ gives
\[
 T_0=\frac{10^6\,235^4 C_{234}}{500^{234}},\qquad
 10^{29}235^4 C_{234}<156\cdot500^{234}.
\]
The displayed integer inequality is exactly
$T_0<156/10^{23}$.  In particular,
\begin{equation}
 \frac{10}{9}T_0<\frac{1560}{9\cdot10^{23}}
       <\frac2{10^{21}},
 \label{supp:first-sum-bound}
\end{equation}
where the last comparison reduces to $1560<1800$.

Let $\rho_{\max}=\max_{1\leq i\leq4}\rho_i(234)$, and define three
exact rational upper bounds
\[
 A_{\rm coarse}=\frac{10}{9}T_0,\qquad
 A_{\rm common}=\frac{T_0}{1-\rho_{\max}},\qquad
 A_{\rm component}=\sum_{i=1}^4
                \frac{t_i(234)}{1-\rho_i(234)}.
\]
Decreasing ratios imply
$t_i(234+6j)\leq t_i(234)\rho_i(234)^j$.  Summing these geometric
majorants and using \eqref{supp:initial-ratio-bound} proves
\begin{equation}
 \mathfrak a\leq A_{\rm component}
 \leq A_{\rm common}<A_{\rm coarse}<\frac2{10^{21}}.
 \label{supp:geometric-bounds}
\end{equation}
The second inequality holds because
$1-\rho_i(234)\geq1-\rho_{\max}>0$.  The common-ratio estimate keeps
only the largest initial ratio, whereas the componentwise estimate retains
each one; both refine the coarse factor $10/9$.

\paragraph{An exact partial sum and a positive remainder bound.}
There are exactly $62$ allowed indices from $234$ through $600$, because
$600=234+6\cdot61$.  Define the finite rational sum
\begin{equation}
 P_{600}=\sum_{i=1}^4\sum_{j=0}^{61}t_i(234+6j)
 \label{supp:partial-600}
\end{equation}
and the positive rational number
\begin{equation}
 R_{606}=\sum_{i=1}^4\frac{t_i(606)}{1-\rho_i(606)},\qquad
 \rho_i(606)=\left(\frac{613}{607}\right)^4
             \left(\frac{102}{101}\right)^{p_i}\beta_i^6.
 \label{supp:remainder-606}
\end{equation}
Every numerator is positive and
$0<\rho_i(606)<\rho_i(234)<1/10$, so $R_{606}>0$.  Starting the
geometric estimate at the next mode $606$, rather than at the last included
mode $600$, gives
\begin{equation}
 P_{600}<\mathfrak a\leq P_{600}+R_{606}.
 \label{supp:partial-tail-bound}
\end{equation}
The first inequality is strict because all subsequent summands are
positive.  The number $R_{606}$ is an upper bound for the remainder, not
an assertion that the remainder equals it.

This refinement is no larger than the componentwise bound.  For the
$i$-th component put $r=\rho_i(234)$.  Then
$t_i(606)\leq t_i(234)r^{62}$ and $\rho_i(606)\leq r$, while each term
of its partial sum is at most $t_i(234)r^j$.  Consequently
\[
 \sum_{j=0}^{61}t_i(234+6j)
       +\frac{t_i(606)}{1-\rho_i(606)}
 \leq t_i(234)\left(\sum_{j=0}^{61}r^j+
                         \frac{r^{62}}{1-r}\right)
 =\frac{t_i(234)}{1-r}.
\]
Summing over $i$ places $P_{600}+R_{606}$ below
$A_{\rm component}$ in \eqref{supp:geometric-bounds}.  These finite bounds
are exact rationals; a rounded display that makes
$P_{600}$ and $P_{600}+R_{606}$ look equal does not change
\eqref{supp:partial-tail-bound}.

\paragraph{Rational propagation to the radial errors.}
The refined bounds can be substituted directly into the propagation
estimates of Section~\ref{tail:propagation}.  They do not replace its
analytic mode bounds or the finite premises
\eqref{tail:finite-core-premise}: the retained second-order $B_1$ norm
is still bounded by $6000$, and the weighted third-order $D_i$ norm by
$2\cdot10^6$.  Recall also that
\begin{equation}
 W_*=4\sum_{k=234,240,\ldots}w_k
 \leq4\sum_{k=234,240,\ldots}(k+1)^4w_k=4\mathfrak a,
 \label{supp:W-bound}
\end{equation}
so replacing $W_*$ by $4\mathfrak a$ is an upper bound, not an identity.
Collecting the integral and mode-norm contributions in
\eqref{tail:f-estimate} and using $\pi<22/7$ gives the rational
expression, for $A,W\geq0$,
\begin{align}
 \mathcal F(A,W)={}&2\frac{22}{7}(3W+16A)\notag\\
 &+\frac12\bigl(60000W+8A+12(2\cdot10^6)W
       +48(6000+W)A\notag\\
 &\hspace{30mm}+8(2\cdot10^6)A+16A^2\bigr).
 \label{supp:rational-error-expression}
\end{align}
The terms are respectively the source integral, the $S_2$ error, the
linear $D_i$ error, the retained and omitted mixed $D_iH_{k_i}$ errors,
and the quadratic $D_iD_j$ error.  The bound $\pi<22/7$ and
\eqref{tail:f-estimate} give
$|f_0-\widetilde f_0|_{2,t}\leq\mathcal F(\mathfrak a,W_*)$.
Every coefficient of $\mathcal F$ is nonnegative.  Thus any one of the
rational upper bounds $A$ above, for which $\mathfrak a\leq A$, satisfies
\[
 |f_0-\widetilde f_0|_{2,t}
 \leq\mathcal F(\mathfrak a,W_*)\leq\mathcal F(A,4A).
\]
Direct collection of the rational terms gives the exact identity
\begin{equation}
 \mathcal F(A,4A)=A\bigl(56\,264\,180+104A\bigr).
 \label{supp:rational-error-coefficient}
\end{equation}
For example, the source term contributes $176A$; the remaining linear
terms contribute $56\,264\,004A$, and the quadratic terms contribute
$104A^2$.

Put $A_0=2/10^{21}$.  Each rational upper bound for $\mathfrak a$ above is
less than $A_0$, and the coefficient in
\eqref{supp:rational-error-coefficient} increases with $A$.
The exact comparison
\[
 56\,264\,180+\frac{208}{10^{21}}
 <56\,264\,181<6\cdot10^7
\]
therefore supplies the claimed coefficient bound.  More explicitly,
\[
 \mathcal F(A_0,4A_0)
 =\frac{112\,528\,360}{10^{21}}+\frac{416}{10^{42}}
 <\frac{12}{10^{14}}.
\]
The order-two norm is
$|F|_{2,t}=|F(t)|+|F'(t)|+|F''(t)|/2$.  Hence the second derivative,
unlike the value and first derivative, requires the factor $2$:
\[
 |f_0''-\widetilde f_0''|
 \leq2|f_0-\widetilde f_0|_{2,t}
 <\frac{24}{10^{14}}=2.4\cdot10^{-13}.
\]
For the angular function, \eqref{tail:e-estimate} similarly yields the
entirely rational comparison
\[
 |e''-\widetilde e''|
 \leq2\bigl(16(2\cdot10^6)A_0+32A_0^2\bigr)
 =\frac{128\,000\,000}{10^{21}}+\frac{256}{10^{42}}
 <\frac{13}{10^{14}}=1.3\cdot10^{-13}.
\]
The value and first derivative are bounded by their norms without this
factor.  Thus these rational alternatives preserve the error threshold
$10^{-11}$ for all six quantities in \eqref{tail:uniform-error} on $[b,l]$.
For $t\geq l$, the exact continuation \eqref{tail:exterior-formula}, rather
than a uniform absolute error at unbounded radius, remains the applicable
statement.
